\section{Extended limitations}
\label{S-sec:s-limitations}

The chemical analogy guides computational design without modeling a physical process. The vacancy has Langmuir form, but its capacities were fixed in advance rather than derived from adsorption equilibria. It is one of several soft occupancy rules; maximal marginal relevance on the matches performs as well in the comparison. Full retention helps but is not optimal in every synthetic setting. The release rule is optimal only for its worst case over surrogate trade-offs; the null behind its shape channel is not calibrated on real scans, and a weighted sum with one weight fitted on training data releases 0.5\% more cells (\cref{S-sec:s-verification}).

Verification uses the observed point clouds, making surface-using methods the appropriate comparison group; every frontend in \main{\cref*{tab:main}} receives the same verification.

With a handful of inliers, NGR outperforms the baselines in the specified synthetic tests (69\% against at most 59\% on a synthetic benchmark frozen before this study; 29 against 4 of 100 fresh inputs with three inliers among 4000 matches; \cref{S-sec:s-synthetic}), and our Python MAC adaptation obtains part of this gain once it uses NGR's count readout. On real pairs with at most ten true inliers, all methods have low success rates. No frontend consistently leads on real scans thinned to 5--20 inliers.

In its frozen implementation the \ngr{} frontend is among the slowest in our CPU runs. For \ls{}, restricting growth to the initial supports and compiling the code gives a frontend time of 0.30\,s at sixteen threads on 3DLoMatch+FPFH and a timed pipeline cost of 0.39\,s, while preserving the frozen outcomes on all 15\,035 pairs. Learned correspondences have less local products, and the compiled frontend still takes 1.0--3.7\,s at four threads. With common verification, \ls{} has the second-highest total count over nine cohorts, behind SC$^2$-PCR.

\paragraph{Protocol and scope.} The \ngr{} frontend and finishing step were fixed before the protocol. Their development used public-benchmark evaluations, so those benchmarks were not held out for these components (\cref{S-sec:s-protocol}). The two components added here, the finite site capacity and the release assay, were each pre-registered in writing and decided only on a training cohort of 3000 FPFH pairs from 82 3DMatch training scenes, disjoint from every test scene. Each was then run once on the test cohorts.

The optional learned assay was only partly covered by the protocol: its 26 candidate features were listed before the protocol was in place, so it is reported separately and used for no training-free claim. The comparisons of the diversity and release rules were added after the main evaluation; each was pre-registered before it was run, and neither changes the reported pipeline. One further pre-registered component, filling the finishing step's unused ICP slots, passed its margin on the training cohort, but its single test run raised the ceiling by 58 of 12\,035 pairs and the released count by 2, so the reported pipeline does not use it. The cube advantage of the frozen synthetic benchmark does not reproduce with fresh seeds (62 against 63 for MAC with its original parameters), and on 3DMatch thinned to five true inliers among 500 matches SC$^2$-PCR registers significantly more pairs than \ls{} (84 against 63, $p=0.007$; \cref{S-sec:s-synthetic}). The main paper's MAC results use the authors' C++ code, which has no released ETH setting; our earlier Python adaptation, which exceeds the 16\,GiB/600\,s limit on 19--95\% of the pairs in four of the five learned-correspondence cohorts, remains in parts of this supplement (\cref{S-sec:s-maccpp}). With eight test scenes absolute recall is uncertain, although the paired gains are stable (\cref{S-sec:s-results}). Published numbers use different evaluation conditions (\cref{S-sec:s-literature}), and the timings compare specific implementations (\cref{S-sec:s-cost}).
